\documentclass[aps,prmaterials,onecolumn,superscriptaddress,floatfix]{revtex4-2}

\usepackage{amsmath,amssymb}
\usepackage{booktabs}
\usepackage{xcolor}
\usepackage[colorlinks=true,
            linkcolor=blue!60!black,
            citecolor=green!50!black,
            urlcolor=blue!70!black]{hyperref}

% Auto-generated by scripts/prm_make_result_assets.py; do not edit.
\newcommand{\PRMSelectedVariant}{\texttt{g111}}
\newcommand{\PRMSelectedValidationMAE}{0.371}
\newcommand{\PRMSelectedValidationMAELow}{0.357}
\newcommand{\PRMSelectedValidationMAEHigh}{0.385}
\newcommand{\PRMSelectedTestMAE}{0.382}
\newcommand{\PRMSelectedTestMAELow}{0.377}
\newcommand{\PRMSelectedTestMAEHigh}{0.387}
\newcommand{\PRMIdCvDARTMAE}{0.418}
\newcommand{\PRMIdCvSchNetMAE}{0.571}
\newcommand{\PRMIdCvDescriptorMAE}{0.625}
\newcommand{\PRMIdCvDescriptorName}{LightGBM}
\newcommand{\PRMIdCvDARTLowEnergyMAE}{0.440}
\newcommand{\PRMIdCvDARTLowEnergyRecall}{87.1}
\newcommand{\PRMPairCvDARTMAE}{0.450}
\newcommand{\PRMPairCvSchNetMAE}{0.628}
\newcommand{\PRMPairCvDescriptorMAE}{0.674}
\newcommand{\PRMPairCvDescriptorName}{LightGBM}
\newcommand{\PRMPairCvDARTLowEnergyMAE}{0.480}
\newcommand{\PRMPairCvDARTLowEnergyRecall}{86.0}
\newcommand{\PRMHostCvDARTMAE}{0.938}
\newcommand{\PRMHostCvSchNetMAE}{6.703}
\newcommand{\PRMHostCvDescriptorMAE}{1.273}
\newcommand{\PRMHostCvDescriptorName}{Histogram GB}
\newcommand{\PRMHostCvDARTLowEnergyMAE}{1.284}
\newcommand{\PRMHostCvDARTLowEnergyRecall}{69.2}
\newcommand{\PRMDopantCvDARTMAE}{0.844}
\newcommand{\PRMDopantCvSchNetMAE}{1.692}
\newcommand{\PRMDopantCvDescriptorMAE}{1.331}
\newcommand{\PRMDopantCvDescriptorName}{Histogram GB}
\newcommand{\PRMDopantCvDARTLowEnergyMAE}{1.041}
\newcommand{\PRMDopantCvDARTLowEnergyRecall}{71.6}
\newcommand{\PRMChemistryBlockDARTMAE}{0.386}
\newcommand{\PRMChemistryBlockSchNetMAE}{0.508}
\newcommand{\PRMChemistryBlockDescriptorMAE}{0.642}
\newcommand{\PRMChemistryBlockDescriptorName}{LightGBM}
\newcommand{\PRMChemistryBlockDARTLowEnergyMAE}{0.305}
\newcommand{\PRMChemistryBlockDARTLowEnergyRecall}{75.0}
\newcommand{\PRMUQTestMAE}{0.377}
\newcommand{\PRMUQNLL}{1.128}
\newcommand{\PRMUQCRPS}{0.324}
\newcommand{\PRMUQErrorSpearman}{0.546}
\newcommand{\PRMUQAURC}{0.177}
\newcommand{\PRMUQExcessAURC}{0.095}
\newcommand{\PRMUQFiftyCoverage}{52.0}
\newcommand{\PRMUQFiftyWidth}{0.336}
\newcommand{\PRMUQEightyCoverage}{78.3}
\newcommand{\PRMUQEightyWidth}{0.752}
\newcommand{\PRMUQNinetyCoverage}{87.2}
\newcommand{\PRMUQNinetyWidth}{1.141}
\newcommand{\PRMUQNinetyFiveCoverage}{93.7}
\newcommand{\PRMUQNinetyFiveWidth}{1.847}
\newcommand{\PRMPreferenceAccuracy}{91.0}
\newcommand{\PRMPreferenceMarginMAE}{0.749}
\newcommand{\PRMPreferenceMarginSpearman}{0.954}
\newcommand{\PRMGlobalRegret}{0.139}
\newcommand{\PRMGlobalExactAccuracy}{69.5}
\newcommand{\PRMWithinTypeExactAccuracy}{73.6}
\newcommand{\PRMWithinTypeTopTwoAccuracy}{91.4}
\newcommand{\PRMWithinTypeRegret}{0.103}


\begin{document}

\title{Supplemental Material for ``Defect-aware graph learning of impurity
incorporation energetics and chemical transferability in two-dimensional
materials''}

\author{Yiming Hua}
\affiliation{Wuhan University, Wuhan, China}

\date{\today}

\maketitle

\section{Post-hoc SchNet graph-readout sensitivity}

The prespecified periodic SchNet comparator follows the original
energy-oriented construction~\cite{schutt2018}: its learned atomwise scalar
contributions are summed to produce the graph prediction.  Defect formation
energy is instead localized and non-extensive with respect to supercell atom
count, and prior atomistic studies show that graph pooling can materially
affect localized or intensive quantum-property prediction
~\cite{buterez2023pooling}.  We therefore performed one bounded post-hoc
sensitivity analysis after completing the main comparison.

The intervention changed only the graph readout from atomwise addition to mean
pooling.  It retained all five immutable host-held-out folds, seeds 342--344,
the SchNet interaction stack, learned atomic-number input, optimizer, learning
rate schedule, host-balanced sampler, target normalization, target-noise
stream, gradient clipping, and 150-epoch budget.  Predictions were averaged
over the three seeds within each fold.  The pooled mean-minus-add difference in
absolute error was evaluated with 20,000 bootstrap samples clustered by the 44
host materials.

% Auto-generated; do not edit.
\begin{table*}[t]
\caption{Post-hoc SchNet graph-readout sensitivity under the fixed
host-held-out protocol. Each fold prediction averages seeds 342--344. The
difference is mean-readout MAE minus additive-readout MAE, in eV. The pooled
host-cluster bootstrap 95\% interval for this difference is [-7.704, -1.256]~eV.
The prespecified additive-readout result remains the main comparator.}
\label{tab:schnet-readout}
\centering
\begin{tabular}{lrrrr}
\toprule
Partition & $n$ & Add MAE & Mean MAE & Mean $-$ add \\
\midrule
Fold 1 & 2023 & 3.882 & 2.179 & -1.704 \\
Fold 2 & 2019 & 8.199 & 2.068 & -6.132 \\
Fold 3 & 2069 & 5.484 & 2.530 & -2.954 \\
Fold 4 & 2057 & 5.466 & 2.468 & -2.998 \\
Fold 5 & 2056 & 10.474 & 3.501 & -6.973 \\
\textbf{Pooled OOF} & 10224 & \textbf{6.703} & \textbf{2.552} & \textbf{-4.151} \\
\bottomrule
\end{tabular}
\end{table*}


Mean pooling reduced the host-held-out MAE in all
\PRMSchNetMeanBetterFolds{}/5 folds.  Across individual samples, the Spearman
association between atom count and absolute error changed from
\PRMSchNetAddSampleSizeErrorSpearman{} for additive pooling to
\PRMSchNetMeanSampleSizeErrorSpearman{} for mean pooling.  Across the 44 hosts,
the corresponding association between median atom count and host MAE changed
from \PRMSchNetAddHostSizeErrorSpearman{} to
\PRMSchNetMeanHostSizeErrorSpearman{}.  These descriptive associations are
consistent with a size-sensitive additive readout, but atom count covaries with
host chemistry and therefore does not identify a standalone causal mechanism.
The intervention establishes that readout choice contributes to the extreme
additive-SchNet error; it neither replaces the prespecified main comparator nor
attributes the remaining DART--SchNet difference to a single component.

\bibliography{references}

\end{document}
